\subsection{The standard resolution}

In this section, one supposes that the ring A is an (associative and unital) algebra over a commutative ring \(k\) . For \(n \geqslant 0\) we denote by \(\mathbf{B}_n\) the tensor product over \(k\) of \((n + 2)\) modules equal to A. We regard it as an (A, A)-bimodule by endowing it with the left (resp. right) A-module structure induced by the left (resp. right) A-module structure of the first (resp. last) factor of the tensor product.

For \(n \geqslant 1\) , we define bimodule homomorphisms \(d_n^i\) (for \(0 \leqslant i \leqslant n\) ) and \(d_n\) from \(B_n\) to \(B_{n-1}\) , by the formulas:

\[
d _ {n} ^ {i} (x _ {0} \otimes \ldots \otimes x _ {n + 1}) = x _ {0} \otimes \ldots \otimes x _ {i} x _ {i + 1} \otimes \ldots \otimes x _ {n + 1}, \quad 0 \leqslant i \leqslant n,
\]

\[
d _ {n} = \sum_ {i = 0} ^ {n} (- 1) ^ {i} d _ {n} ^ {i}.
\]

It is clear that

\[
d _ {n - 1} ^ {i} \circ d _ {n} ^ {j} = d _ {n - 1} ^ {j - 1} \circ d _ {n} ^ {i} \quad \text { for } \quad i <   j
\]

and consequently

\[
d _ {n - 1} \circ d _ {n} = \sum_ {0 \leqslant i <   j \leqslant n} (- 1) ^ {i + j} d _ {n - 1} ^ {i} \circ d _ {n} ^ {j} + \sum_ {0 \leqslant j \leqslant i \leqslant n - 1} (- 1) ^ {i + j} d _ {n - 1} ^ {i} \circ d _ {n} ^ {j} = 0.
\]

Therefore, if we set \(\mathbf{B}_n = 0\) for \(n < 0\) and \(d_n = 0\) for \(n \leqslant 0\) , the sequence \((\mathbf{B}_n, d_n)\) defines a complex of (A, A)-bimodules (X, p.~43), which will be denoted \(\mathbf{B}(\mathbf{A})\) . For every left A-module \(\mathbf{M}\) we denote by \(\mathbf{B}(\mathbf{A}, \mathbf{M})\) the complex formed by the \(\mathbf{B}_n \otimes_{\mathbf{A}} \mathbf{M}\) and the \(d_n \otimes 1_{\mathbf{M}}\) , in other words the tensor product complex \(\mathbf{B}(\mathbf{A}) \otimes_{\mathbf{A}} \mathbf{M}_{*}\) ; it is a complex of left A-modules.

We define an A-linear map \(\varepsilon_{\mathbf{M}}\) from \(\mathbf{B}_{0}(\mathbf{A},\mathbf{M}) = (\mathbf{A}\otimes_{k}\mathbf{A})\otimes_{\mathbf{A}}\mathbf{M}\) to \(\mathbf{M}\) by the formula \(\varepsilon_{\mathbf{M}}(a\otimes b\otimes m) = abm\) for \(a,b\in \mathbf{A},m\in \mathbf{M}\) . We have \(\varepsilon_{\mathbf{M}}\circ d_1 = 0\) , so that the graded homomorphism \(\overline{\varepsilon}_{\mathbf{M}}:\mathbf{B}(\mathbf{A},\mathbf{M})\to \mathbf{M}\) , which coincides with \(\varepsilon_{\mathbf{M}}\) in degree 0, is a morphism of complexes of A-modules.

PROPOSITION 12. --- The map \(\overline{\varepsilon}_{\mathbf{M}}: \mathbf{B}(\mathbf{A}, \mathbf{M}) \to \mathbf{M}\) is a homotopism of complexes of \(k\) -modules. In particular, the complex \(\mathbf{B}(\mathbf{A}, \mathbf{M})\) is split over \(k\) , and \((\mathbf{B}(\mathbf{A}, \mathbf{M}), \overline{\varepsilon}_{\mathbf{M}})\) is a left resolution of the A-module M.

For \(n \geqslant 0\) , let us define a \(k\) -linear map \(s_n: \mathbf{B}_n \to \mathbf{B}_{n+1}\) by the formula:

\[
s _ {n} \left(x _ {0} \otimes \dots \otimes x _ {n + 1}\right) = 1 \otimes x _ {0} \otimes \dots \otimes x _ {n + 1} \quad \text { for } \quad x _ {0}, \dots , x _ {n + 1} \in A.
\]

It is a homomorphism of right A-modules, which satisfies the identities:

\[
d _ {n + 1} ^ {i} \circ s _ {n} = s _ {n - 1} \circ d _ {n} ^ {i - 1} \quad \mathrm{for} \quad n \geqslant 1, \quad 1 \leqslant i \leqslant n + 1,
\]

\[
d _ {n + 1} ^ {0} \circ s _ {n} = 1 _ {\mathrm{B} _ {n}} \quad \text { for } \quad n \geqslant 1,
\]

and consequently we have

\[
d _ {n + 1} \circ s _ {n} + s _ {n - 1} \circ d _ {n} = 1 _ {\mathrm{B} _ {n}} \quad \text { for } \quad n \geqslant 1.\tag{16}
\]

Moreover, we have

\[
d _ {1} \circ s _ {0} (x _ {0} \otimes x _ {1}) = x _ {0} \otimes x _ {1} - 1 \otimes x _ {0} x _ {1} \quad \text { for } \quad x _ {0}, x _ {1} \in A.\tag{17}
\]

Let us denote by \(\eta: \mathbf{A} \to \mathbf{A} \otimes_k \mathbf{A}\) the map defined by \(\eta(a) = 1 \otimes a\) , and \(\overline{\eta}: \mathbf{A} \to \mathbf{B}(\mathbf{A})\) the morphism of complexes which coincides with \(\eta\) in degree 0. It is clear that \(\overline{\varepsilon}_{\mathbf{A}} \circ \overline{\eta} = 1_{\mathbf{A}}\) ; formulas (16) and (17) show that \(d \circ s + s \circ d = 1_{\mathrm{B(A)}} - \overline{\eta} \circ \overline{\varepsilon}_{\mathbf{A}}\) . Setting \(\overline{\eta}_{\mathrm{M}} = \overline{\eta} \otimes 1_{\mathrm{M}}\) , \(d_{\mathrm{M}} = d \otimes 1_{\mathrm{M}}\) and \(s_{\mathrm{M}} = s \otimes 1_{\mathrm{M}}\) , we deduce that \(\overline{\varepsilon}_{\mathrm{M}} \circ \overline{\eta}_{\mathrm{M}} = 1_{\mathrm{M}}\) and \(d_{\mathrm{M}} \circ s_{\mathrm{M}} + s_{\mathrm{M}} \circ d_{\mathrm{M}} = 1_{\mathrm{B(A,M)}} - \overline{\eta}_{\mathrm{M}} \circ \overline{\varepsilon}_{\mathrm{M}}\) . In other words, (X, p.~33, def. 5), \(\overline{\varepsilon}_{\mathrm{M}}\) is a homotopism of complexes of \(k\) -modules. The other assertions of the proposition follow immediately.

DEFINITION 3. --- The left resolution (B(A, M), \(\overline{\varepsilon}_{\mathrm{M}}\) ) of M is called the standard resolution of the A-module M.

If A and M are projective (resp. free, resp. flat) k-modules, the standard resolution B(A, M) is a projective (resp. free, resp. flat) resolution of M.

